\section{Experimental Setup}

\subsection{Input requirements}
DELT-Hit uses three main input categories:

(1) Library definition files (.xlsx): These files specify building block structures in SMILES format, reaction SMIRKS templates, DNA constant sequences, and barcode-to-building block mapping tables. They are typically provided as Excel workbooks with standardized sheet names for automated parsing.

(2) Raw sequencing data (fastq.gz): FASTQ files generated by Illumina or compatible sequencing platforms are supported in compressed and uncompressed formats.

(3) Experimental metadata (.yaml): These files define selection conditions, target proteins, control experiments, and replicate groupings for downstream statistical analysis and quality control.


\subsection{Library architecture considerations}
DELT-Hit supports several library architectures:

Single-display libraries: Linear synthetic schemes in which DNA tags are appended after each reaction cycle.

Dual-display libraries: Architectures in which compounds are displayed on both complementary DNA strands.

Multi-cycle libraries: Schemes with multiple synthetic steps, branching reactions, and multiple building block incorporation sites.

The framework validates library architecture consistency and issues warnings for potential problems such as incomplete reaction definitions or missing building blocks.


\subsection{Library chemistry and reaction definition}
Reaction cycles used during library construction are specified with reaction SMIRKS in the Excel configuration file. DELT-Hit translates these definitions into an explicit reaction graph that supports arbitrary reaction sequences, branching pathways, and multi-step transformations.

DELT-Hit accommodates libraries of varying complexity, including architectures with multiple reaction cycles, provided that each synthetic route resolves to a single, well-defined final product. Libraries with non-standard reaction schemes, unconventional building block representations, or complex precursor handling may require additional customization and validation.

Library enumeration depends on accurate SMIRKS definitions. DELT-Hit provides templates for commonly used DEL reactions, but library-specific adaptations may require advanced cheminformatics expertise to ensure chemically correct product formation and reliable structure reconstruction.


\subsection{Quality control parameters}
DELT-Hit monitors quality metrics throughout the workflow:

Demultiplexing efficiency: Percentage of sequencing reads assigned to valid barcode combinations, typically \textgreater{}70\% for high-quality datasets.

Barcode recovery rates: Coverage of individual barcodes to identify imbalances or failed assignments.

Statistical significance assessment: Multiple-testing correction and false discovery rate control using established RNA-seq-derived methods.

Replicate consistency: Correlation analysis across biological and technical replicates.

Chemical reaction validation: Verification that reaction definitions produce the expected products.
